% Dinâmica: 2 blocos (e 1 força)
% Faro, seg 15 dez 2025 21:20:53 
% Written by Tordar 
% pstricks2pdf.sh 2blocos
\documentclass{article}
\pagestyle{empty}

\usepackage{pst-all} 
\usepackage{pstricks-add}
\usepackage{pst-slpe}
\usepackage{graphicx}
\input{apnum}

% Bloco
\newcommand{\bloco}[5][]{% #1=options, #2=cx, #3=cy, #4=totalWidth, #5=height
  % Left frame (half of total width, full height)
  \psframe[linewidth=1pt,linecolor=gray,fillstyle=slope,slopebegin=white,slopeend=gray,slopesteps=15]
    (!#2 #4 2 div sub #3 #5 2 div sub)                     % lower-left of total area
    (!#2 #3 #5 2 div add)                                  % upper-right of left frame (center x, top)  
  % Right frame (half of total width, full height)
  \psframe[linewidth=1pt,linecolor=gray,fillstyle=slope,slopebegin=gray,slopeend=white,slopesteps=15]
    (!#2 #3 #5 2 div sub)                                  % lower-left of right frame (center x, bottom)
    (!#2 #4 2 div add #3 #5 2 div add)                     % upper-right of total area
  \psframe[#1](!#2 #4 2 div sub #3 #5 2 div sub)(!#2 #4 2 div add #3 #5 2 div add) 
}

\begin{document}

% Altura e largura do bloco inferior: 
\evaldef\wbi{5}
\evaldef\hbi{2}
% Altura e largura do bloco superior: 
\evaldef\wbii{3}
\evaldef\hbii{1.5}
% Coordenadas do bloco inferior: 
\evaldef\xbi{4}
\evaldef\ybi{2+\hbi/2}
\evaldef\xmi{\xbi+\hbi+1}
% Coordenadas do bloco superior: 
\evaldef\xbii{\xbi}
\evaldef\ybii{\ybi+\hbi/2+\hbii/2}
\evaldef\xmii{\xbii+\hbii+1}
% Coordenadas da força:
\evaldef\xfe{\xbi-\wbi/2}
\evaldef\xfs{\xfe-2}
\evaldef\xf{(\xfs+\xfe)/2}
\evaldef\yf{\ybi+0.5}
% Círculo do zoom:
\evaldef\xz{\xbi}
\evaldef\yz{\ybi+0.5*\hbi}
%\evaldef\xf{(\xfs+\xfe)/2}
%\evaldef\yf{\ybi+0.5}
% Aceleração:
%\evaldef\xas{\xbii+0.6*\wbii}
%\evaldef\xae{\xas+1.5}
%\evaldef\xa{(\xas+\xae)/2}
%\evaldef\ya{\ybii+0.5}

% Let's go:
\begin{pspicture}(0,0)(10,10)
% Chão: 
\psframe[linewidth=0pt,linecolor=lightgray,dimen=inner,fillstyle=hlines*,fillcolor=lightgray,hatchangle=45](0,0)(10,2) 
% Força: 
\psline[linewidth=3pt,linecolor=blue]{->}(\xfs,\ybi)(\xfe,\ybi) 
% Blocos: 
\bloco[linewidth=2pt,linecolor=black]{\xbi}{\ybi}{\wbi}{\hbi}
\bloco[linewidth=2pt,linecolor=black]{\xbii}{\ybii}{\wbii}{\hbii}
% Zoom: 
\pscircle(\xz,\yz){0.2}
% Clipping:
\begin{psclip}{\pscircle[linewidth=2pt](10,4){2}}
\pspolygon[linewidth=2pt,linecolor=lightgray,dimen=inner,fillstyle=hlines*,fillcolor=lightgray,hatchangle=45](8,6)(12,6)(12,5)(11,4)(10,5.5)(8.5,4.2)(8,5)
\pspolygon[linewidth=2pt,linecolor=lightgray,dimen=inner,fillstyle=hlines*,fillcolor=lightgray,hatchangle=45](8,2)(12,2)(12,4)(11.5,4.5)(9.5,4)(8.5,4.5)(8,4)
\end{psclip}
% Par de forças de atrito: 
\psline[linewidth=3pt,linecolor=blue]{->}(11,5)(13,5)\rput{0}(12,5.5){\scalebox{1.5}{$\mathbf{F}_{at}$}}
\psline[linewidth=3pt,linecolor=blue]{<-}(9,2.5)(11,2.5)\rput*{0}(10,3){\scalebox{1.5}{$-\mathbf{F}_{at}$}}
% Arco do clip ao zoom: 
\pscurve[linewidth=2pt,linecolor=black]{->}(8.3,5)(6,7)(4,4.2)
% Texto: 
\rput{0}(\xmi,\ybi){\scalebox{1.5}{$m_2$}}
\rput{0}(\xmii,\ybii){\scalebox{1.5}{$m_1$}}
\rput{0}(\xf,\yf){\scalebox{1.5}{$\mathbf{F}$}}
\end{pspicture}

\end{document}
